\documentclass[a4paper]{article}
\usepackage{amsfonts}
\usepackage{amsmath}
\usepackage{amssymb}
\usepackage{graphicx}     

\newcommand{\dint}{\displaystyle\int}
\newcommand{\llim}{\lim\limits}

\begin{document}
  
\noindent \large Auteur: Oubayy Ahale          \\
\noindent \large Studierichting: Informatica   \\
\noindent \large Oefeningenreeks 3C nr. 18.19  \\

\medskip

\normalsize

$ \llim_{x \rightarrow +\infty} \left( x - x^2 \ln \left( 1 + \dfrac{1}{x} \right) \right) $\\

$ \Leftrightarrow \llim_{x \rightarrow +\infty} \left( 1 - 2x \ln \left( 1 + \dfrac{1}{x} \right) + x^2 \cdot \dfrac{1}{1 + \dfrac{1}{x}} \cdot \dfrac{x \cdot 0 - 1 \cdot 1}{x^2} \right) $\\

$ \Leftrightarrow \llim_{x \rightarrow +\infty} \left( 1 - 2x \ln \left( 1 + \dfrac{1}{x} \right) + \dfrac{1}{1 + \dfrac{1}{x}} \right) $\\

Stel $ t = \dfrac{1}{x} \Rightarrow  x = \dfrac{1}{t} $\\

Wanneer $ x \to \infty $, geldt dat $ t \to 0 $\\

We herschrijven de limiet in termen van $ t $:\\

$ \llim_{t \rightarrow 0} \left( \dfrac{1}{t} - \dfrac{1}{t^2} \ln (1 + t) \right) $\\

$ = \llim_{t \rightarrow 0} \left( \dfrac{t - \ln (1+t)}{t^2} \right) $\\

$ = \llim_{t \rightarrow 0} \dfrac{1 - \dfrac{1}{1 + t}}{2t} $\\

$ = \llim_{t \rightarrow 0} \dfrac{t}{2t + 2t^2} $\\

$ = \llim_{t \rightarrow 0} \dfrac{1}{2 + 2t} $\\

$ = \dfrac{1}{2} $\\

\begin{thebibliography}{999}
%\bibitem{a} Referentie
\end{thebibliography}
\end{document}